\documentclass[11pt,a4paper]{article}
\usepackage[margin=25mm]{geometry}
\usepackage[T1]{fontenc}
\usepackage{lmodern,microtype,amsmath,amssymb,amsthm,booktabs,array}
\usepackage[hidelinks]{hyperref}
\usepackage{url}
\newtheorem{theorem}{Theorem}
\newtheorem{lemma}[theorem]{Lemma}
\newtheorem{proposition}[theorem]{Proposition}
\theoremstyle{definition}\newtheorem{definition}[theorem]{Definition}
\newcommand{\HH}{\mathrm H}
\newcommand{\PP}{\mathrm P}
\newcommand{\DD}{\mathrm D}
\newcommand{\Res}{\operatorname{Res}}
\newcommand{\nnz}{\operatorname{nnz}}
\setlength{\emergencystretch}{1.5em}
\title{Term-minimal realizations of nine line-type configurations\protect\\in quadratic mass-action systems}
\author{Tatsuhiko Takimoto}
\date{24 October 2026 \quad Draft v1}
\begin{document}
\maketitle
\begin{abstract}
Consider two-species mass-action systems with reactant molecularity at most two,
unrestricted products, and a positive equilibrium whose Jacobian is the zero
matrix. We count nonzero monomials separately in the two component equations,
in concentration coordinates. After moving the equilibrium to the origin for
geometric analysis, we require the two homogeneous quadratic components to have
nonzero resultant. For nine configurations of invariant-line types, the minimum
term counts are $8,8,8,8,8,8,8,6,6$. Rational witnesses attain every bound.
The lower bounds follow from a nonnegative six-parameter representation and an
exhaustion of the eight support faces with fewer than eight terms. Relative to
the examples in Yeung and Erban's classification paper, including its six-term
example, five upper bounds decrease and four existing upper bounds are proved
optimal within this domain. The count concerns ODE terms; minimum reaction
counts are a separate problem.
\end{abstract}

\section{The question and the result}
Yeung and Erban~\cite{YE} identify nine realizable geometric configurations at
positive equilibria with zero linearization in planar second-order chemical
systems. Their discussion asks whether the nonzero-term counts of the displayed
examples can be reduced. This paper gives exact minima for the nine
invariant-line configurations under the explicit nonzero-resultant condition
below. Table~\ref{tab:minima} states the result and compares the examples.

The geometric background is the classification attributed to
Lyagina~\cite{Lyagina} and developed in Newton's invariant-line
description~\cite{Newton}. We use the line-type notation explained in
Section~\ref{sec:lines}, and distinguish configurations by their multisets of
line types. This avoids depending on inconsistent numerical labels in the
source examples. The proof below does not require a new proof of the historical
classification or an identification of its entire domain with ours.

\begin{table}[ht]
\centering
\small
\begin{tabular}{@{}lrrrrl@{}}
\toprule
Configuration & Row$^{a}$ & Table 1 & Best$^{b}$ & Minimum & Comparison \\
\midrule
$\HH$ & 1 & 9 & 9 & 8 & decreases \\
$\PP_2\HH$ & 4 & 8 & 8 & 8 & attained before \\
$\DD,\HH$ & 5 & 10 & 10 & 8 & decreases \\
$\DD,\PP_2\HH$ & 8 & 8 & 8 & 8 & attained before \\
$\PP_1,\PP_2,\HH$ & 11 & 9 & 9 & 8 & decreases \\
$\PP_1,\HH,\HH$ & 12 & 9 & 9 & 8 & decreases \\
$\PP_1,\PP_2,\PP_2\HH$ & 14 & 8 & 8 & 8 & attained before \\
$\PP_1,\HH,\PP_2\HH$ & 15 & 8 & 6 & 6 & attained before \\
$\PP_1,\PP_2\HH,\PP_2\HH$ & 16 & 8 & 8 & 6 & decreases \\
\bottomrule
\end{tabular}
\caption{Minimum numbers of component monomials. $^{a}$Row label in Table 1
of~\cite{YE}, used only to locate the example. $^{b}$Best count among that
table and the six-term example in its Section 4. The configurations here are
computed from the coefficients, not copied from the printed labels.}
\label{tab:minima}
\end{table}

\section{Chemical domain and the quantity minimized}
Write $\dot x=f(x,y)$ and $\dot y=g(x,y)$, with real coefficients and total
degree at most two. A reaction
$\alpha X+\beta Y\longrightarrow\gamma X+\delta Y$ has a positive rate
constant and $\alpha+\beta\leq2$. We impose no bound on
$\gamma+\delta$. Consequently the coefficient of a monomial not containing
$x$ must be nonnegative in $f$, and the coefficient of a monomial not
containing $y$ must be nonnegative in $g$.
These sign conditions also suffice: realize each nonzero coefficient by a
reaction that adds or removes one molecule of the corresponding species.
A negative coefficient has that species in its reactant, so removal is
permitted. This constructs a network with positive rate constants.

Let $(x_*,y_*)$ be a selected positive equilibrium and require
$Df(x_*,y_*)=Dg(x_*,y_*)=0$. Positive changes of concentration units place it
at $(1,1)$. Such changes multiply each coefficient by a positive number and
preserve the number of nonzero terms. With $u=x-1$ and $v=y-1$, the shifted
equations have the form
\begin{equation}\label{eq:shifted}
 \dot u=P(u,v)=Au^2+Buv+Cv^2,\qquad
 \dot v=Q(u,v)=Du^2+Euv+Fv^2.
\end{equation}
Translation is used to study invariant lines; all term counts remain in
the concentration variables $x,y$.

In the monomial order $(1,x,y,x^2,xy,y^2)$, the coefficient vectors are
\begin{align}
 a&=(A+B+C,-2A-B,-B-2C,A,B,C),\nonumber\\
 b&=(D+E+F,-2D-E,-E-2F,D,E,F).\label{eq:original}
\end{align}
The chemical sign conditions give the following parameterization:
\begin{equation}\label{eq:cone}
\begin{split}
 \theta&=(s,t,c,q,r,d)\in\mathbb R_{\geq0}^{6},\\
 a&=(s,-2s-t,t,s+t+c,-t-2c,c),\\
 b&=(q,r,-2q-r,d,-r-2d,q+r+d),\\
 (A,B,C,D,E,F)&=(s+t+c,-t-2c,c,d,-r-2d,q+r+d).
\end{split}
\end{equation}
Every chemical quadratic system with the specified equilibrium and zero
Jacobian has this form, since $s=a_{00}$, $t=a_{01}$, $c=a_{02}$,
$q=b_{00}$, $r=b_{10}$ and $d=b_{20}$. Conversely, direct substitution in
\eqref{eq:cone} verifies all these properties.

\begin{definition}\label{def:domain}
A parameter vector in~\eqref{eq:cone} is admissible if the homogeneous
resultant $\Res(P,Q)$ is nonzero. Its term count is
\[
 M(\theta)=\nnz(a)+\nnz(b).
\]
For a specified line-type multiset $T$, its minimum is taken over all
admissible real parameter vectors with type $T$.
\end{definition}
The resultant condition means that $P$ and $Q$ have no common zero on the
complex projective line; it excludes zero components and common factors.
It is part of the domain, not a conclusion inferred from an isolated real
equilibrium. The same monomial appearing in both components contributes two
to $M$. Positive concentration and time scalings preserve $M$; arbitrary
affine changes of concentration variables are not operations in this
minimization problem. A reaction may contribute to both components, so
$M$ does not count reactions.

\section{Invariant-line types}\label{sec:lines}
A finite slope $z=v/u$ defines an invariant line precisely when
\begin{equation}\label{eq:h}
 h(z)=zP(1,z)-Q(1,z)=Cz^3+(B-F)z^2+(A-E)z-D=0.
\end{equation}
At any real root, $p(z)=P(1,z)\ne0$ by the resultant condition. Set
\begin{equation}\label{eq:rho}
 n(z)=E+(2F-B)z-2Cz^2,\qquad
 \rho(z)=\frac{n(z)}{p(z)}=1-\frac{h'(z)}{p(z)}.
\end{equation}
For a simple root, the four symbols in Table~\ref{tab:types} record the
range of $\rho$. A double root has $\rho=1$ and is denoted $\DD$.
The multiplicity condition is essential: $\DD$ does not include triple
roots. All roots and the vertical direction are counted as projective
directions, so each line appears once.

\begin{table}[ht]
\centering
\begin{tabular}{@{}lllll@{}}
\toprule
Symbol & $\PP_1$ & $\PP_2$ & $\HH$ & $\PP_2\HH$\\
\midrule
Condition (simple root) & $\rho>1$ & $0<\rho<1$ & $\rho<0$ & $\rho=0$\\
\bottomrule
\end{tabular}
\caption{Operational definitions of the four simple line types.}
\label{tab:types}
\end{table}

These quantities have a direct geometric meaning. Near a finite invariant
line $v=zu$, put $w=v-zu$. Linearizing the scalar trajectory equation gives
$dw/du=\rho w/u$, whose first variation is proportional to
$|u|^\rho$. This distinguishes approach tangent to the line at the origin
($\PP_1$), relative approach at large distance ($\PP_2$), vanishing
transverse distance at large distance ($\HH$), and the boundary case
($\PP_2\HH$). The definitions used in the proof are the algebraic
conditions of~\eqref{eq:rho}, rather than sector counts alone.

The vertical line $u=0$ is invariant if and only if $C=0$. In this case
$F>0$ and $B\leq0$, so it is a simple projective root and
\begin{equation}\label{eq:infinity}
 \rho_\infty=B/F.
\end{equation}
Indeed the swapped slope polynomial is
$wQ(w,1)-P(w,1)$, whose derivative at $w=0$ is $F-B>0$.
We use unordered multisets of these symbols. No assertion that all nine
are distinct under arbitrary local topological conjugacy is needed.

\section{Exact upper bounds}
\begin{theorem}\label{thm:main}
In the domain of Definition~\ref{def:domain}, the minima for the nine
configurations in Table~\ref{tab:minima}, in the displayed order, are
\[
 8,8,8,8,8,8,8,6,6.
\]
\end{theorem}
Table~\ref{tab:witness} gives integer parameter vectors that attain these
counts. Together with~\eqref{eq:cone}, the table specifies every coefficient
of every witness exactly. Positive common scaling of the vector amounts
to a change of time units.

\begin{table}[ht]
\centering\small
\begin{tabular}{@{}clrrrrrrr@{}}
\toprule
ID & Configuration & $s$ & $t$ & $c$ & $q$ & $r$ & $d$ & $M$\\
\midrule
1 & $\HH$ & 0&0&3&1&0&2&8\\
2 & $\PP_2\HH$ & 1&0&0&2&0&1&8\\
3 & $\DD,\HH$ & 0&0&9&4&0&2&8\\
4 & $\DD,\PP_2\HH$ & 2&0&0&3&0&1&8\\
5 & $\PP_1,\PP_2,\HH$ & 0&0&56&25&0&12&8\\
6 & $\PP_1,\HH,\HH$ & 1&1&0&0&0&2&8\\
7 & $\PP_1,\PP_2,\PP_2\HH$ & 31&0&0&33&0&64&8\\
8 & $\PP_1,\HH,\PP_2\HH$ & 1&0&0&0&0&1&6\\
9 & $\PP_1,\PP_2\HH,\PP_2\HH$ & 1&0&0&1&0&0&6\\
\bottomrule
\end{tabular}
\caption{Exact witnesses, with parameters in~\eqref{eq:cone}.}
\label{tab:witness}
\end{table}

\clearpage
For completeness, Table~\ref{tab:h} records their slope polynomials and
transverse ratios. In row 5 the ratios follow the root order
$2,2/7,3/8$. In rows 6--8 the paired signs have the same order in the
roots and the ratios. The only real root in row 1 is $2$, and row 2 has
no real finite root. In rows 3 and 4, $1/3$ and $1/2$, respectively, are
double roots.

\begin{table}[ht]
\centering\small
\renewcommand{\arraystretch}{1.2}
\begin{tabular}{@{}clll@{}}
\toprule
ID & $h(z)$ & Finite roots : $\rho$ & $\rho_\infty$\\
\midrule
1 & $(z-2)(3z^2-3z+1)$ & $2:-4/3$ & ---\\
2 & $-3z^2+3z-1$ & none & $0$\\
3 & $(z-2)(3z-1)^2$ & $2:-16/9;\quad 1/3:1$ & ---\\
4 & $-(2z-1)^2$ & $1/2:1$ & $0$\\
5 & $(z-2)(7z-2)(8z-3)$ & $2,2/7,3/8:-25/14,7/10,48/35$ & ---\\
6 & $-3z^2+6z-2$ & $1\pm1/\sqrt3:4\pm3\sqrt3$ & $-1/2$\\
7 & $-97z^2+159z-64$ & $(159\pm\sqrt{449})/194:1\pm\sqrt{449}/31$ & $0$\\
8 & $-z^2+3z-1$ & $(3\pm\sqrt5)/2:1\pm\sqrt5$ & $0$\\
9 & $z-z^2$ & $0:0;\quad 1:2$ & $0$\\
\bottomrule
\end{tabular}
\caption{Algebraic checks of the line types. A dash means that the vertical
direction is not invariant.}
\label{tab:h}
\end{table}

The homogeneous resultants, in the same order, are
\[
 9,\ 9,\ 1296,\ 64,\ 1960000,\ 4,\ 9042049,\ 1,\ 1.
\]
All are nonzero. Thus each displayed system is admissible; substituting its
parameters into~\eqref{eq:cone} gives the asserted count and chemical
realization. In row 6, direct substitution in~\eqref{eq:rho} is also a
convenient way to check the two finite types.

\section{Lower bounds over real parameters}\label{sec:lower}
The lower bound does not depend on searching a finite grid of coefficients.
It follows from the support of the nonnegative coordinates in~\eqref{eq:cone}.
Encode the positive entries of $(s,t,c)$ by a mask with weights $(1,2,4)$.
Use the same weights for $(q,r,d)$. Positivity and the absence of cancellation
in the sums in~\eqref{eq:cone} give the component counts
\begin{equation}\label{eq:masks}
\begin{array}{c|rrrrrrr}
 \text{mask}&1&2&3&4&5&6&7\\\hline
 \text{component terms}&3&4&5&3&5&5&6.
\end{array}
\end{equation}
A zero mask would make one component identically zero and is excluded.
It follows at once that $M\geq6$.

\begin{lemma}\label{lem:low}
An admissible system with $M<8$ has configuration either
$\{\PP_1,\HH,\PP_2\HH\}$ or
$\{\PP_1,\PP_2\HH,\PP_2\HH\}$.
\end{lemma}
\begin{proof}
Equation~\eqref{eq:masks} leaves just eight pairs of masks:
\[
 (1,1),(1,4),(4,1),(4,4),(1,2),(2,1),(4,2),(2,4).
\]
Three pairs have both $s=q=0$, namely $(4,4),(4,2),(2,4)$; in these pairs,
one has $P(1,1)=s=0$ and $Q(1,1)=q=0$. The components have a common
projective zero, contradicting admissibility.

Interchanging the species exchanges the masks and preserves the line types.
For a nonzero finite slope this follows by conjugating the trajectory slope
map $f(z)=Q(1,z)/P(1,z)$ with $z\mapsto1/z$; the derivative at its fixed
point remains $\rho$. The horizontal and vertical charts give the same
statement for the remaining directions. Hence only $(1,1),(1,4),(1,2)$
need be considered. Write their positive coordinates as $a,b>0$.

For $(1,1)$, $P=au^2$ and $Q=bv^2$. The directions are $0,a/b,\infty$
and the ratios are $0,2,0$. The configuration is
$\{\PP_1,\PP_2\HH,\PP_2\HH\}$.

For $(1,4)$, $P=au^2$ and $Q=b(u-v)^2$. Here
\[
 h(z)=-bz^2+(a+2b)z-b,\qquad \Delta=a^2+4ab>a^2.
\]
The finite roots have ratios $1\pm\sqrt\Delta/a$ and the vertical
direction has ratio zero. The configuration is
$\{\PP_1,\HH,\PP_2\HH\}$.

For $(1,2)$, $P=au^2$ and $Q=bv(v-u)$. The directions
$0,(a+b)/b,\infty$ have ratios $-b/a,2+b/a,0$, again giving
$\{\PP_1,\HH,\PP_2\HH\}$. In all three families the resultant is
$a^2b^2>0$. This exhausts all eight faces for all real choices of their
positive coordinates.
\end{proof}

\begin{proof}[Proof of Theorem~\ref{thm:main}]
The last two configurations attain six terms by Table~\ref{tab:witness},
and~\eqref{eq:masks} proves that fewer are impossible. Lemma~\ref{lem:low}
excludes fewer than eight terms for each of the other seven configurations.
Their eight-term witnesses complete the proof.
\end{proof}

\section{Computation and assurance}\label{sec:assurance}
The calculation that led to the witnesses decomposes the parameter cone into
$7\cdot7=49$ nonempty support faces and considers the nine target types on
each face. Polynomial solving supplies candidates. The minima depend only
on the nine rational witnesses and the lower-bound argument in
Section~\ref{sec:lower}; unresolved feasibility cases on larger faces are
not used in that argument.

For a rational witness, an exact certificate contains the component
coefficients, a factorization of $h$, root multiplicities, rational isolating
intervals, and signs of $p,n,n-p$ at each real root. A checker using rational
arithmetic verifies the factorization, squarefreeness of its factors,
Sturm root counts, interval signs or exact zero identities, the vertical
direction, and the resultant. It also checks the equilibrium, zero Jacobian,
chemical signs, term count, and a reaction realization. Irrational roots of
rational polynomials require isolating intervals, not floating-point
classification. Lower-bound certificates encode the component support
counts and the three symbolic families in Lemma~\ref{lem:low}.

The assurance consists of these exact certificate checks and the displayed
algebraic proof; no proof-assistant formalization or human peer
review is claimed. Provenance, checker versions, and the bounded literature
search are recorded separately in the accompanying
\texttt{VERIFICATION-RECORD.md}.

\section{Independent verification}
The nine minima were also obtained in a separate verification that used the
problem specification and its stated definitions, while withholding the
research-side minima, proofs, implementations, and certificates. It constructed
integer-coefficient witnesses and derived a lower bound by classifying all
$2^6=64$ zero-or-positive supports of the nonnegative parameters. Removing
supports with a zero component leaves the 49 faces considered above. The eight
supports with $M<8$ reduce to the same three positive parameter families after
excluding common factors and interchanging the species. All nine independently
obtained minima agree with Theorem~\ref{thm:main}.

The separate derivation was checked with a symbolic-algebra implementation
using SymPy and a separately written Python standard-library implementation.
The reported results agree. These checks were produced within one verification;
their agreement is not an additional independent human review.

The shared starting point includes the admissible domain, the nonnegative
six-parameter representation~\eqref{eq:cone}, its component-term counting rule,
and the invariant-line tests~\eqref{eq:h}--\eqref{eq:infinity}. From these
supplied preparations, the verification separately determined the minima,
constructed witnesses, and proved the lower bounds. The specification also
included published examples, so the verification was not blind to those known
examples. Neither derivation includes a full-text comparison of the
classification domains in Newton's and Lyagina's original papers.

\section{Comparison and scope}
The comparison in Table~\ref{tab:minima} uses coefficients rather than class
numbers. In Table 1 of~\cite{YE}, the printed line-type labels in rows
14--16 differ from the types obtained from those rows' coefficients using
\eqref{eq:h}--\eqref{eq:infinity}. Section 4 of that paper gives
\[
 \dot x=(x-y)^2,\qquad \dot y=(y-1)^2.
\]
Its six terms give configuration $\{\PP_1,\HH,\PP_2\HH\}$, as can also
be seen by exchanging the species in row 8 of Table~\ref{tab:witness}.
The text calls this class 12. We retain the explicit system and its
computed configuration as the comparison, without proposing a corrected
historical numbering convention. Four parabolic sectors and two hyperbolic
sectors alone do not distinguish the line-type multisets being minimized.

Lemma 6(ii) of~\cite{YE} contains a further boundary discrepancy: for
$\eta_2>1$ it prints $\eta_1<1/\eta_2$, whereas its proof, equation (3.43),
gives $\eta_1\leq1/\eta_2$. The admissible system
$P=\tfrac12u^2$, $Q=(u-v)^2$ has finite slopes $1/2,2$ and attains equality.
Its resultant is $1/4$ and its ratios are $-2,4,0$, including the vertical
direction. The nonnegative parameter cone in~\eqref{eq:cone} includes this
valid boundary; no strict version of that source inequality is used here.

Related minimality problems use different constraints. Chen and
Tang~\cite{CT} classify minimal-dimensional degenerate zero-one networks;
their two-dimensional condition concerns the stoichiometric subspace and
their degeneracy condition differs from a zero planar Jacobian at a
selected equilibrium. Plesa and Sprott~\cite{PS} study small monomial
counts for three-dimensional chemical systems with chaos. Beldiman and
Halanay~\cite{BH} classify minimal canard networks; their canard-point
conditions require nonzero cross derivatives. These results do not supply
the nine minima of Table~\ref{tab:minima} in the material examined.

The result covers the real coefficient domain of Definition~\ref{def:domain}.
It leaves systems with zero resultant, equilibria on the boundary of the
concentration quadrant, other molecularity constraints, and minimum numbers
of reactions outside its scope. The bounded literature review has not
located the same nine minima. Full-text review of the historical sources
and a broader priority search remain incomplete.

\paragraph{Use of AI.}
AI tools assisted the algebraic search, certificate software, literature
comparison, and drafting. The mathematical claims are supported by the
proof and exact data described above. Responsibility for their review and
any decision to publish rests with the author.

\paragraph{Acknowledgment of prior work.}
The question comes from Yeung and Erban, and the geometric framework from
the work of Lyagina and Newton cited above. The author is grateful for
these foundations. This acknowledgment concerns their published work and
does not imply their involvement in or endorsement of this draft.

\begin{thebibliography}{9}
\bibitem{YE} N. Yeung and R. Erban,
\emph{Classification of strongly non-hyperbolic critical points of planar
second-order chemical reaction systems}, arXiv:2610.06130v1 (2026).
\url{https://arxiv.org/abs/2610.06130v1}.

\bibitem{Lyagina} L. S. Lyagina,
\emph{The integral curves of the equation
$y'=(ax^2+bxy+cy^2)/(dx^2+exy+fy^2)$},
Uspekhi Matematicheskikh Nauk \textbf{6}, no. 2(42), 171--183 (1951),
in Russian. \url{https://www.mathnet.ru/eng/rm6840}.

\bibitem{Newton} T. A. Newton,
\emph{Two Dimensional Homogeneous Quadratic Differential Systems},
SIAM Review \textbf{20}(1), 120--138 (1978).
\url{https://doi.org/10.1137/1020007}.

\bibitem{CT} Y. Chen and X. Tang,
\emph{Characterization of Minimal Degenerate Zero-One Reaction Networks},
arXiv:2608.18509v1 (2026). \url{https://arxiv.org/abs/2608.18509v1}.

\bibitem{PS} T. Plesa and J. C. Sprott,
\emph{Simple chemical systems with chaos}, arXiv:2601.02787v1 (2026).
\url{https://arxiv.org/abs/2601.02787v1}.

\bibitem{BH} M.-M. Beldiman and A.-D. Halanay,
\emph{Canard points in planar mass-action networks: a flux criterion,
a duality, and a complete classification}, arXiv:2610.27188v1 (2026).
\url{https://arxiv.org/abs/2610.27188v1}.
\end{thebibliography}
\end{document}
